% !TeX root = ../../Book2.tex
\section{Bessel 方程及其标准解}
\label{b2:sec:A201}

对径向函数作 Fourier 变换时，角向积分不再是一个指数函数，而会产生随半径振荡的特殊函数。同一函数也出现在圆盘振动的分离变量法中：角向模态固定后，径向方程多出一个反平方势项。Bessel 方程正是连接这两种计算的常微分方程。

\ProseParagraphBreak
先看圆盘上的一个实际计算。设 \(\kappa>0\)，在半径为 \(a>0\) 的圆盘内考虑
\(-\Delta u=\kappa^2u\)。写
\(x=\rho\cos\theta\)、\(y=\rho\sin\theta\)，其中 \(0<\rho<a\)。
链式法则给出
\[
 \partial_x=\cos\theta\,\partial_\rho
 -\frac{\sin\theta}{\rho}\partial_\theta,\quad
 \partial_y=\sin\theta\,\partial_\rho
 +\frac{\cos\theta}{\rho}\partial_\theta.
\]
将两个一阶算子各自作用两次并相加，混合导数项抵消，
而对系数求导留下 \(\rho^{-1}\partial_\rho\)，所以
\[
 \Delta=\partial_\rho^2+\rho^{-1}\partial_\rho
 +\rho^{-2}\partial_\theta^2.
\]
固定整数角向模态 \(m\in\mathbb Z\)，取
\(u(\rho,\theta)=R(\rho)e^{\mathrm im\theta}\)；
整数条件保证 \(\theta\) 增加 \(2\pi\) 后函数不变。
由于 \(\partial_\theta^2e^{\mathrm im\theta}=-m^2e^{\mathrm im\theta}\)，原方程化为
\begin{equation}\label{b2:eq:A2-disk-radial}
 \rho^2R''(\rho)+\rho R'(\rho)
 +(\kappa^2\rho^2-m^2)R(\rho)=0.
\end{equation}
再令 \(r=\kappa\rho\)、\(y(r)=R(r/\kappa)\)，便得到
\[
 r^2y''(r)+ry'(r)+(r^2-\nu^2)y(r)=0,
 \quad \nu=|m|.
\]
模态 \(m=0\) 是圆盘上的径向振动，此时 \(\nu=0\)；
其他模态产生 \(\nu^2/r^2\) 项。这个推导同时说明，原点处的适当条件来自圆盘中心的正则性，
不是在奇点上任意给定两个常微分方程初值。

\subsection{级数解与原点行为}
\label{b2:subsec:A20101}

取实阶参数 \(\nu\in\mathbb R\)。对 \(r>0\) 和 \(C^2\) 函数 \(y:(0,\infty)\to\mathbb C\)，\term[bessel-fangcheng]{Bessel 方程}[Bessel equation]为
\[r^2y''+ry'+(r^2-\nu^2)y=0.\]
本附录主要处理实阶参数；先取 \(\nu>-1\)。\chapref{b2:ch:06}已经定义正实参数的 Gamma 积分。分部积分给出 \(\Gamma(s+1)=s\Gamma(s)\)、\(\Gamma(1)=1\)，Gauss 积分给出 \(\Gamma(1/2)=\sqrt\pi\)。这些等式足以处理下面全部实阶级数与几何常数。

\ProseParagraphBreak
在原点附近，暂时略去较低影响的 \(r^2y\) 项，试代 \(y=r^\sigma\)，
得到指标方程 \(\sigma^2-\nu^2=0\)。
为了构造其中的 \(r^\nu\) 分支，进一步寻找
\(y=r^\nu\sum_{k\geq0}c_kr^{2k}\)。
采用偶次步长，是因为方程中的乘法项 \(r^2y\) 恰好把幂次提高二。
算子 \(r^2\partial_r^2+r\partial_r-\nu^2\) 作用于
\(r^{2k+\nu}\) 时给出
\[
 \bigl((2k+\nu)^2-\nu^2\bigr)r^{2k+\nu}
 =4k(k+\nu)r^{2k+\nu}.
\]
于是 \(k=0\) 项自动消失，其他项要求
\[
 c_k=-\frac{c_{k-1}}{4k(k+\nu)}
 =\frac{(-1)^kc_0}
 {4^kk!\,(\nu+1)\cdots(\nu+k)}\quad(k\geq1).
\]
条件 \(\nu>-1\) 保证这些分母不为零。
取 \(c_0=2^{-\nu}/\Gamma(\nu+1)\)，再把连乘积写成 Gamma 比值，
就得到下面的标准归一化。

\Needspace{6\baselineskip}
\begin{definition}[第一类 Bessel 函数]
\label{b2:def:A2-bessel-j}
对 \(\nu>-1\)、\(r>0\)，定义
\[
 J_\nu(r)=\sum_{k=0}^\infty
 \frac{(-1)^k}{k!\,\Gamma(k+\nu+1)}\left(\frac r2\right)^{2k+\nu}.
\]
称它为 \term[diyilei-bessel-hanshu]{第一类 Bessel 函数}[Bessel function of the first kind]。实幂在正半轴取通常正值。
\end{definition}
\symidx{\(J_\nu(r)\)：第一类 Bessel 函数}

\Needspace{6\baselineskip}
\begin{proposition}[级数与微分关系]
\label{b2:prop:A2-bessel-series}
设 \(\nu>-1\)、\(r>0\)，\(J_\nu\) 按\cref{b2:def:A2-bessel-j}的正实幂级数定义。该级数在正半轴每个紧区间上连同任意阶导数一致收敛，并满足 Bessel 方程。原点附近有
\[J_\nu(r)=\frac{(r/2)^\nu}{\Gamma(\nu+1)}(1+O(r^2)).\]
此外，对全部 \(\nu>-1\)、\(r>0\)，有
\[
 \begin{aligned}
 \frac{\mathrm d}{\mathrm dr}(r^{-\nu}J_\nu(r))
 &=-r^{-\nu}J_{\nu+1}(r),\\
 \frac{\mathrm d}{\mathrm dr}(r^{\nu+1}J_{\nu+1}(r))
 &=r^{\nu+1}J_\nu(r).
 \end{aligned}
\]
\end{proposition}
\begin{proof}
相邻项绝对值之比为
\(r^2/[4(k+1)(k+\nu+1)]\)，在任意有界 \(r\) 区间上一致趋零。
更准确地，\(r^{-\nu}J_\nu(r)\) 是关于 \(r^2\) 的收敛半径为无穷的幂级数。
在 \(0<c\leq r\leq C<\infty\) 上，任意固定次微分只增加关于 \(k\) 的多项式因子，
再乘一个由 \(c,C,\nu\) 控制的因子，所以各导数级数也一致收敛。
这既保证逐项微分，也说明结论没有把可能奇异的端点 \(r=0\) 包含进来。
将第 \(k\) 项记为 \(c_kr^{2k+\nu}\)，代入方程后，
\(r^{2k+\nu}\) 的系数恰为
\(4k(k+\nu)c_k+c_{k-1}\)，由 Gamma 递推为零。
首项非零，下一项与首项之比为 \(-r^2/[4(\nu+1)]\)，所以还有较具体的展开
\[
 J_\nu(r)=\frac{(r/2)^\nu}{\Gamma(\nu+1)}
 \left(1-\frac{r^2}{4(\nu+1)}+O(r^4)\right).
\]

第一条导数关系中，\(k=0\) 项求导为零。把其余项的指标写成 \(k=\ell+1\)，得到
\[
 \begin{aligned}
 \frac{\mathrm d}{\mathrm dr}(r^{-\nu}J_\nu(r))
 &=\sum_{k=1}^\infty
 \frac{(-1)^k\,2k\,r^{2k-1}}
 {2^{2k+\nu}k!\,\Gamma(k+\nu+1)}\\
 &=-\sum_{\ell=0}^\infty
 \frac{(-1)^\ell r^{2\ell+1}}
 {2^{2\ell+\nu+1}\ell!\,\Gamma(\ell+\nu+2)}
 =-r^{-\nu}J_{\nu+1}(r).
 \end{aligned}
\]
第二条关系不需变动指标，而是约去 Gamma 递推中的因子：
\[
 \begin{aligned}
 \frac{\mathrm d}{\mathrm dr}(r^{\nu+1}J_{\nu+1}(r))
 &=\sum_{k=0}^\infty
 \frac{(-1)^k(2k+2\nu+2)r^{2k+2\nu+1}}
 {2^{2k+\nu+1}k!\,\Gamma(k+\nu+2)}\\
 &=\sum_{k=0}^\infty
 \frac{(-1)^kr^{2k+2\nu+1}}
 {2^{2k+\nu}k!\,\Gamma(k+\nu+1)}
 =r^{\nu+1}J_\nu(r).
 \end{aligned}
\]
\end{proof}

例如 \(J_0(r)=1-r^2/4+r^4/64-\cdots\)、\(J_1(r)=r/2-r^3/16+\cdots\)。
一般实阶时，光滑的是加权函数 \(r^{-\nu}J_\nu(r)\) 在原点的偶次展开；
不能因此断言 \(J_\nu\) 本身在原点具有任意阶导数。
特别地，\(-1<\nu<0\) 时首项已经发散。
两条导数关系也可改写为
\[
 J_\nu'(r)=\frac{\nu}{r}J_\nu(r)-J_{\nu+1}(r),\quad
 J_{\nu+1}'(r)=J_\nu(r)-\frac{\nu+1}{r}J_{\nu+1}(r).
\]
这些关系把相邻阶数连接起来，但没有改变各公式中 \(\nu>-1\)、\(r>0\) 的范围。

\ProseParagraphBreak
半整数阶给出可以直接检验常数的例子。由
\[
 \Gamma(k+1/2)=\frac{(2k)!}{4^kk!}\sqrt\pi
\]
逐项比较余弦和正弦级数，得到
\[
 \begin{aligned}
 J_{-1/2}(r)
 &=\sqrt{\frac2{\pi r}}
 \sum_{k=0}^\infty\frac{(-1)^kr^{2k}}{(2k)!}
 =\sqrt{\frac2{\pi r}}\cos r,\\
 J_{1/2}(r)
 &=\sqrt{\frac2{\pi r}}
 \sum_{k=0}^\infty\frac{(-1)^kr^{2k+1}}{(2k+1)!}
 =\sqrt{\frac2{\pi r}}\sin r.
 \end{aligned}
\]
它们解释了三维径向 Fourier 公式为何最终只含正弦。一旦阶数不是半整数，保留 Bessel 记号通常比反复展开级数更清楚。

\ProseParagraphBreak

对于圆盘模态 \(\nu=|m|\)，原点有界的这一支与中心正则性相容：
\(J_{|m|}(\kappa\rho)e^{\mathrm im\theta}\) 是
\(\rho^{|m|}e^{\mathrm im\theta}\) 乘以关于 \(\rho^2\) 的收敛级数，
而前者在 \(m\geq0\) 时等于 \((x+\mathrm iy)^m\)，在 \(m<0\) 时等于
\((x-\mathrm iy)^{|m|}\)。
因此这一模态确实光滑地穿过圆心；\(m=0\) 时级数还给出径向导数在中心为零。
方程在 \(r=0\) 有奇点，不能像普通初值问题那样任意指定 \(y(0),y'(0)\)。
在每个离开零点的区间上，二阶常微分方程仍具有通常的两维解空间；
另一支将在大参数归一化后构造，届时才能说明原点有界性为何排除它。

\subsection{积分表示}
\label{b2:subsec:A20102}

对实数 \(a,b>0\)，记 \(B(a,b)=\int_0^1u^{a-1}(1-u)^{b-1}\,\mathrm du\)，称为\term[beta-hanshu]{Beta 函数}[beta function]。两端可积恰好使用了参数正性。
\Needspace{6\baselineskip}
\begin{lemma}[实参数 Beta 积分]\label{b2:lem:A2-beta-integral}
设实数 \(a,b>0\)，
\(B(a,b)=\int_0^1t^{a-1}(1-t)^{b-1}\,\mathrm dt\)，
并对 \(s>0\) 定义 \(\Gamma(s)=\int_0^\infty t^{s-1}e^{-t}\,\mathrm dt\)。
则 \(B(a,b)=\Gamma(a)\Gamma(b)/\Gamma(a+b)\)。
\end{lemma}
\begin{proof}
在第一象限的二重 Gamma 积分中作 \(x=su,y=s(1-u)\)，其中 \(s>0,0<u<1\)，Jacobi 因子为 \(s\)。因此
\[
 \begin{aligned}
 \Gamma(a)\Gamma(b)
 &=\int_0^\infty\int_0^1
 e^{-s}s^{a+b-1}u^{a-1}(1-u)^{b-1}\,\mathrm du\,\mathrm ds\\
 &=\Gamma(a+b)B(a,b).
 \end{aligned}
\]
非负性允许 Tonelli 换序，两个 Gamma 积分有限且 \(\Gamma(a+b)>0\)，相除即得结论。
\end{proof}

\Needspace{6\baselineskip}
\begin{proposition}[Poisson 积分表示]
\label{b2:prop:A2-poisson-bessel}
设实数 \(\nu>-1/2\)、\(r>0\)，\(J_\nu\) 为级数归一化的第一类 Bessel 函数，则
\[
 J_\nu(r)=\frac{(r/2)^\nu}{\sqrt\pi\,\Gamma(\nu+1/2)}
       \int_{-1}^1e^{irt}(1-t^2)^{\nu-1/2}\,\mathrm dt.
\]
\end{proposition}
\begin{proof}
先在有限区间上展开指数函数。所有项的绝对值之和被
\(e^r(1-t^2)^{\nu-1/2}\) 支配，而 \(\nu>-1/2\) 保证这个权在两端可积，
所以可以交换求和与积分。奇数项因对称性消失。
偶数项的积分经 \(u=t^2\) 换元，再用\cref{b2:lem:A2-beta-integral}，得到
\[
 \begin{aligned}
 \int_{-1}^1t^{2k}(1-t^2)^{\nu-1/2}\,\mathrm dt
 &=\int_0^1u^{k-1/2}(1-u)^{\nu-1/2}\,\mathrm du\\
 &=\frac{\Gamma(k+1/2)\Gamma(\nu+1/2)}
 {\Gamma(k+\nu+1)}.
 \end{aligned}
\]
将它与指数级数的系数 \((-1)^kr^{2k}/(2k)!\) 相乘，
再用 \(\Gamma(k+1/2)/(2k)!=\sqrt\pi/(4^kk!)\)，得到
\[
 \int_{-1}^1e^{\mathrm irt}(1-t^2)^{\nu-1/2}\,\mathrm dt
 =\sqrt\pi\,\Gamma(\nu+1/2)
 \sum_{k=0}^\infty
 \frac{(-1)^k(r/2)^{2k}}{k!\,\Gamma(k+\nu+1)}.
\]
乘以陈述中的前因子，恰好恢复 \(J_\nu\) 的级数。
\end{proof}

奇数部分的积分为零，所以也可以把指数替换成 \(\cos(rt)\)，从而直接看出
\(J_\nu(r)\) 为实数。去掉振荡，再对权作一次 Beta 积分，得到
\[
 |J_\nu(r)|\leq\frac{(r/2)^\nu}{\Gamma(\nu+1)}
 \quad(\nu>-1/2,\ r>0).
\]
但这个界在大 \(r\) 时丢失了振荡抵消，不能用来判断 Fourier 核的衰减。
参数边界 \(\nu=-1/2\) 也不能直接代入该积分，因为权变成不可积的
\((1-t^2)^{-1}\)；这一阶仍可使用前面已经证明的余弦公式。
要读出固定阶的精确衰减，需要分析端点 \(t=\pm1\)；
它们对应球面上法向与频率方向平行的两个驻点。

\subsection{大参数渐近与第二类解}
\label{b2:subsec:A20103}

以下渐近均固定阶数 \(\nu\)，令正实变量 \(r\to\infty\)；常数可以依赖阶数，不能据此处理 \(\nu\) 与 \(r\) 同时增大的情形。下面的端点计算将衰减幂和复相位同时确定。它不是对一般振荡积分直接套用绝对值界。

\Needspace{6\baselineskip}
\begin{lemma}[幂型端点的首项]\label{b2:lem:A2-endpoint-asymptotic}
设实数 \(a>-1\)，复函数 \(b\) 在 \([0,\infty)\) 光滑且紧支撑，
其中原点处的光滑性指存在向一个开邻域的光滑延拓。
当正实数 \(r\to+\infty\) 时，
\[
 \int_0^\infty e^{-\mathrm irt}t^ab(t)\,\mathrm dt
 =\Gamma(a+1)e^{-\mathrm i\pi(a+1)/2}b(0)r^{-a-1}
 +O(r^{-a-2}),
\]
余项常数可以依赖固定的 \(a,b\)。
\end{lemma}
\begin{proof}
先计算带指数衰减的模型。设
\(F(r)=\int_0^\infty t^ae^{-(1+\mathrm ir)t}\,\mathrm dt\)。
微分及分部积分分别给出
\[
 \begin{aligned}
 F'(r)&=-\mathrm i\int_0^\infty
 t^{a+1}e^{-(1+\mathrm ir)t}\,\mathrm dt,\\
 (1+\mathrm ir)\int_0^\infty
 t^{a+1}e^{-(1+\mathrm ir)t}\,\mathrm dt&=(a+1)F(r).
 \end{aligned}
\]
故 \(F'(r)=-\mathrm i(a+1)F(r)/(1+\mathrm ir)\)，而
\(F(0)=\Gamma(a+1)\)。解这个实参数微分方程，得到
\[
 \begin{aligned}
 F(r)&=\Gamma(a+1)(1+\mathrm ir)^{-a-1}\\
 &=\Gamma(a+1)e^{-\mathrm i\pi(a+1)/2}r^{-a-1}
 +O(r^{-a-2}),
 \end{aligned}
\]
幂取右半平面的连续辐角。微分时 \(t^{a+1}e^{-t}\) 可积；分部积分时 \(t^{a+1}\) 在原点趋零，所以两步均合法。

写 \(b(t)=b(0)e^{-t}+t c(t)\)。分子在0为零，故 \(c\) 在0光滑；
例如可将差商写成
\(\int_0^1[b'(st)+b(0)e^{-st}]\,\mathrm ds\)，看出各阶导数连续。
在无穷远，\(c\) 及其各阶导数指数衰减。
只需证明余项
\(\int_0^\infty e^{-\mathrm irt}t^{a+1}c(t)\,\mathrm dt=O(r^{-a-2})\)。
远离0的部分是光滑快速衰减函数，可分部积分任意多次。
在0附近先取一个于 \(t\leq2/r\) 支撑的光滑截断，
其直接绝对积分由
\(\int_0^{2/r}t^{a+1}\,\mathrm dt\) 控制，不超过 \(Cr^{-a-2}\)。
其余部分用光滑分割放在
\(t\asymp2^j/r\)、\(j\ge0\) 的环带，直到一个固定小尺度。对每个环带分部积分整数 \(N>a+2\) 次；振幅第 \(N\) 阶导数的积分范数至多
\(C(2^j/r)^{a+2-N}\)，包括截断导数在内。
具体地，若环带尺度为 \(\lambda=2^j/r\)，对 \(t^{a+1}\) 求 \(l\) 次导数产生
\(\lambda^{a+1-l}\)，剩余截断导数至多产生 \(\lambda^{-(N-l)}\)，
而积分区间的长度为 \(O(\lambda)\)，合起来就是所列幂次。
该项因此不超过
\[
Cr^{-N}(2^j/r)^{a+2-N}
=Cr^{-a-2}2^{j(a+2-N)}.
\]
指数严格为负，求几何和给 \(Cr^{-a-2}\)，常数与 \(r\) 无关。各环带截断紧支撑，没有边界项。模型加余项便证明公式；证明还显示常数只用到振幅有限个导数及尾部积分界。
\end{proof}

\Needspace{6\baselineskip}
\begin{theorem}[实轴 Bessel 渐近]
\label{b2:thm:A2-bessel-asymptotic}
设 \(J_\nu\) 为级数归一化的第一类 Bessel 函数。
对固定实数 \(\nu>-1/2\)，令
\(\omega_\nu(r)=r-\pi\nu/2-\pi/4\)。
当正实数 \(r\to+\infty\) 时，
\[
 \begin{aligned}
 J_\nu(r)&=\sqrt{\frac2{\pi r}}\cos\omega_\nu(r)+O(r^{-3/2}),\\
 J_\nu'(r)&=-\sqrt{\frac2{\pi r}}\sin\omega_\nu(r)+O(r^{-3/2}).
 \end{aligned}
\]
对每个固定实数 \(\nu\geq0\)，存在唯一
\(Y_\nu\in C^2((0,\infty);\mathbb R)\)，满足
\(r^2Y_\nu''+rY_\nu'+(r^2-\nu^2)Y_\nu=0\) 以及
\[
 \begin{aligned}
 Y_\nu(r)&=\sqrt{\frac2{\pi r}}\sin\omega_\nu(r)+O(r^{-3/2}),\\
 Y_\nu'(r)&=\sqrt{\frac2{\pi r}}\cos\omega_\nu(r)+O(r^{-3/2}).
 \end{aligned}
\]
所有余项均只断言固定阶数时成立，其常数允许依赖 \(\nu\)。
\end{theorem}
\begin{proof}
写 \(\alpha=\nu-1/2>-1\)，在 Poisson 表示的两个端点附近作对称的光滑分割。
两端之外的振幅光滑且紧支撑于 \((-1,1)\)，反复分部积分给出
\(O(r^{-N})\)，其中 \(N\) 可任意大。
在 \(t=1\) 附近令 \(s=1-t\)，则
\[
 (1-t^2)^\alpha=s^\alpha(2-s)^\alpha,\quad
 e^{\mathrm irt}=e^{\mathrm ir}e^{-\mathrm irs}.
\]
选取等于一的端点截断后，光滑因子在 \(s=0\) 的值为 \(2^\alpha\)。
由\cref{b2:lem:A2-endpoint-asymptotic}，这一端的贡献为
\[
 2^\alpha\Gamma(\alpha+1)r^{-\alpha-1}
 e^{\mathrm i[r-\pi(\alpha+1)/2]}+O(r^{-\alpha-2}).
\]
另一端贡献为其共轭。两端相加得到
\[
 2^{\alpha+1}\Gamma(\alpha+1)r^{-\alpha-1}
 \cos\bigl(r-\pi(\alpha+1)/2\bigr)+O(r^{-\alpha-2}).
\]
选择内部积分的分部积分次数 \(N>\nu+3/2\)，
再乘上 Poisson 前因子
\(r^\nu/[2^\nu\sqrt\pi\,\Gamma(\nu+1/2)]\)。
因为
\[
 \nu-\alpha-1=-\tfrac12,\quad
 2^{\alpha+1-\nu}=\sqrt2,\quad
 \frac{\pi(\alpha+1)}2=\frac{\pi\nu}2+\frac\pi4,
\]
恰好得到 \(J_\nu\) 的余弦主项，端点余项变成 \(O(r^{-3/2})\)。

导数式另行从积分表示求出，不对尚未控制导数的余项作形式微分。
积分对 \(r\) 求导时振幅多出 \(\mathrm it\)，
在右端点取值为 \(\mathrm i\)，在左端点取值为 \(-\mathrm i\)。
因此原来合成 \(2\cos\omega_\nu\) 的两项现在合成
\(-2\sin\omega_\nu\)，同样的端点估计仍有 \(O(r^{-3/2})\) 误差。
前因子求导还产生 \(\nu J_\nu(r)/r=O(r^{-3/2})\)，
故它不改变导数的首项。积分求导合法，因为
\(|t|\cdot(1-t^2)^{\nu-1/2}\) 可积。

为构造第二个解，令 \(w=\sqrt r\,y\)、\(q=\nu^2-1/4\)。方程成为 \(w''+w=q r^{-2}w\)。在充分大的 \([R,\infty)\) 上，考虑
\[
 w(r)=\sqrt{2/\pi}\sin\omega_\nu(r)
 +\int_r^\infty\sin(t-r)\frac q{t^2}w(t)\,\mathrm dt.
\]
自由项选择正弦，是为了确定第二支与 \(J_\nu\) 相差四分之一周期的相位；
积分项补偿方程中尚未包含的 \(q/r^2\) 势。
记积分算子为 \(T\)。在有界连续函数空间 \(C_b([R,\infty))\) 中，
\[
 |T\varphi(r)|\leq
 \int_r^\infty\frac{|q|}{t^2}\cdot|\varphi(t)|\,\mathrm dt
 \leq\frac{|q|}{r}\cdot\|\varphi\|_\infty.
\]
取 \(R>\max\{1,2|q|\}\)，便有 \(\|T\|\leq|q|/R<1/2\)。
收缩原理给出唯一有界连续解 \(w\)，且
\(\|w\|_\infty\leq2\sqrt{2/\pi}\)。
自由项和积分核为实值，所以唯一解也为实值。
对积分项求导时，下端点的正弦为零，因此
\[
 (T\varphi)'(r)
 =-\int_r^\infty\cos(t-r)\frac q{t^2}\varphi(t)\,\mathrm dt.
\]
两项尾积分均为 \(O(r^{-1})\)，故
\[
 \begin{aligned}
 w(r)&=\sqrt{2/\pi}\sin\omega_\nu(r)+O(r^{-1}),\\
 w'(r)&=\sqrt{2/\pi}\cos\omega_\nu(r)+O(r^{-1}).
 \end{aligned}
\]
再求一次导数，连续性和可积尾界给出
\((T\varphi)''=qr^{-2}\varphi-T\varphi\)。
所以 \(w\) 确实满足 \(w''+w=qr^{-2}w\)。
令 \(y=w/\sqrt r\)，并用
\[
 y'(r)=r^{-1/2}w'(r)-\tfrac12r^{-3/2}w(r),
\]
就得到所列 \(y,y'\) 的两项渐近。
取某个 \(r>R\) 上的值与导数作为初值，线性常微分方程的存在唯一性将解延拓到每个
\([\delta,\infty)\)、\(\delta>0\)，因而得到整个正半轴上的解。

若有另一个解满足同样两项渐近，两者之差对应的加权函数 \(v,v'\) 在无穷远趋零。
对等价的一阶系统使用\cref{b2:thm:ode-variation-constants}的变常数公式，从终点 \(L>r\) 向 \(r\) 回溯，得到
\[
 \begin{aligned}
 v(r)={}&v(L)\cos(L-r)-v'(L)\sin(L-r)\\
 &+\int_r^L\sin(t-r)\frac q{t^2}v(t)\,\mathrm dt.
 \end{aligned}
\]
令 \(L\to\infty\)，两项边界值趋零；
由于 \(v\) 在大半轴有界，积分可按 \(t^{-2}\) 支配收敛。
于是得到不含自由项的方程 \(v=Tv\)。
收缩性迫使 \(v\) 在大半轴为零，再由初值唯一性处处为零。
\end{proof}

对实数 \(\nu\geq0\)，将\cref{b2:thm:A2-bessel-asymptotic}中由正弦相位及其导数渐近唯一确定的实 Bessel 解，
称为 \term[dierlei-bessel-hanshu]{第二类 Bessel 函数}[Bessel function of the second kind]，记为 \(Y_\nu\)。
这一归一化确定了正实轴上的第二类解，并与常用约定一致。
定义、积分表示和固定阶渐近的标准记法可对照 NIST《Digital Library of Mathematical Functions》
\cite[\S\S10.2, 10.9, 10.17]{NationalInstituteofStandardsNIST}。
\symidx{\(Y_\nu(r)\)：第二类 Bessel 函数}

\Needspace{6\baselineskip}
\begin{proposition}[Wronski 行列式与原点的第二支]
\label{b2:prop:A2-bessel-wronskian}
令 \(J_\nu\) 取\cref{b2:def:A2-bessel-j}的级数归一化、\(Y_\nu\) 取\cref{b2:thm:A2-bessel-asymptotic}的实轴正弦渐近归一化。对 \(\nu\ge0\)、\(r>0\)，有
\[J_\nu(r)Y_\nu'(r)-J_\nu'(r)Y_\nu(r)=\frac2{\pi r}.\]
因此二者线性无关。原点附近，\(\nu>0\) 时
\[Y_\nu(r)=-\frac{\Gamma(\nu)}\pi\left(\frac2r\right)^\nu(1+o(1)),\]
而 \(Y_0(r)=\frac2\pi\log r+O(1)\)。
\end{proposition}
\begin{proof}
对两解的 Wronski 行列式
\(W=J_\nu Y_\nu'-J_\nu'Y_\nu\) 求导，
交叉项抵消，剩下 \(W'=J_\nu Y_\nu''-J_\nu''Y_\nu\)。
由方程
\(y''=-y'/r-(1-\nu^2/r^2)y\)，可得 \(W'=-W/r\)，
所以 \(rW\) 为常数。用无穷远渐近计算，
\[
 W(r)=\frac2{\pi r}
 \bigl(\cos^2\omega_\nu(r)+\sin^2\omega_\nu(r)\bigr)
 +O(r^{-2}),
\]
故该常数为 \(2/\pi\)，特别地两解线性无关。

由原点级数，存在 \(r_0>0\)，使 \(J_\nu\) 在 \((0,r_0]\) 不为零。
所以在这个区间上
\[
 \left(\frac{Y_\nu}{J_\nu}\right)'=\frac2{\pi rJ_\nu(r)^2}.
\]
令 \(a=2^{-\nu}/\Gamma(\nu+1)>0\)，使用
\(J_\nu(r)=ar^\nu(1+O(r^2))\)，并从 \(r_0\) 积到 \(r\)，得到
\[
 \frac{Y_\nu(r)}{J_\nu(r)}
 =\frac{Y_\nu(r_0)}{J_\nu(r_0)}
 -\frac2{\pi a^2}\int_r^{r_0}s^{-2\nu-1}(1+O(s^2))\,\mathrm ds.
\]
若 \(\nu>0\)，主积分等于
\((r^{-2\nu}-r_0^{-2\nu})/(2\nu)\)。
误差积分由 \(\int_r^{r_0}s^{1-2\nu}\,\mathrm ds\) 控制，因而
\[
 \frac{Y_\nu(r)}{J_\nu(r)}
 =-\frac{r^{-2\nu}}{\pi\nu a^2}
 +\begin{cases}
 O(1),&0<\nu<1,\\
 O(1+|\log r|),&\nu=1,\\
 O(r^{2-2\nu}),&\nu>1.
 \end{cases}
\]
各段余项均为 \(o(r^{-2\nu})\)。
乘回 \(J_\nu\)，再用 \(\Gamma(\nu+1)=\nu\Gamma(\nu)\)，就得到
\[
 -\frac1{\pi\nu a}r^{-\nu}
 =-\frac{\Gamma(\nu)}\pi\left(\frac2r\right)^\nu
\]
作为 \(Y_\nu\) 的首项；中间阶数产生的对数没有改变这一主项。
若 \(\nu=0\)，则 \(a=1\)，被积函数为
\(2/(\pi r)+O(r)\)。
因此 \(Y_0/J_0=(2/\pi)\log r+C+O(r^2)\)。
再乘 \(J_0=1+O(r^2)\)，得到
\(Y_0(r)=(2/\pi)\log r+C+O(r^2|\log r|)\)，特别地有陈述中的有界余项。
\end{proof}

Wronski 常数确定了奇异首项的系数，但仅靠它和原点首项还不能唯一确定第二类解：
\(Y_\nu+cJ_\nu\) 具有同样的 Wronski 行列式和同样的奇异领先项。
无穷远的纯正弦归一化排除了额外的余弦成分，从而固定这个 \(c\)，
也固定了零阶原点展开中的常数 \(C\)。

\ProseParagraphBreak
对 \(\nu\geq0\)，两解组成正半轴上的基本解组。
\(Y_\nu\) 在原点发散，而 \(J_\nu\) 有界，因此任何原点有界的解都是 \(J_\nu\) 的常数倍。
回到\cref{b2:eq:A2-disk-radial}的整数角向模态，这说明中心正则性确实选出前面构造的光滑分支。
无穷远的出入射条件则会选择 \(J_\nu\pm\mathrm iY_\nu\) 的组合，
其首项分别是
\(\sqrt{2/(\pi r)}\,e^{\pm\mathrm i\omega_\nu(r)}\)。
圆盘边值问题与外部散射问题因此不能使用同一条“去掉奇解”的规则。
上述渐近式固定了实阶参数；若阶数与自变量同时增大，还需要控制它们的相对尺度。
